\item \textbf{Barrera de Coulomb para la fusión Deuterio-Tritio.}
Considere la reacción de fusión:
$$ ^{2}_{1}\text{H} + ^{3}_{1}\text{H} \to ^{4}_{2}\text{He} + ^{1}_{0}\text{n} $$
con el núcleo de tritio inicialmente en reposo.
\begin{enumerate}
    \item Asumiendo que los núcleos se fusionan si sus superficies se tocan, determine la distancia entre sus centros al contacto usando el radio aproximado nuclear $R = R_0 A^{1/3}$.
    \item Calcule la energía potencial eléctrica de repulsión culombiana (en $\text{eV}$) a esta distancia.
    \item Si el deuterón se dispara directamente al tritio con la energía justa para vencer la barrera culombiana, calcule la velocidad común de los núcleos al tocarse en términos de la velocidad inicial del deuterón $v_i$.
    \item Use conservación de la energía para determinar la energía cinética mínima inicial del deuterón necesaria para lograr fusión clásica.
    \item Explique por qué en la práctica la fusión termonuclear ocurre a energías promedio significativamente menores que la calculada en (d).
\end{enumerate}